% !TEX root = ../main.tex
%%%%%%%%%%%%%%%%%%%%%%%%%%%%%%%%%%%%%%%%%%%%%%%%%%%%%%%%%%%%%%%%%%%%%%%
% APPENDIX - Validation and reproducibility
%%%%%%%%%%%%%%%%%%%%%%%%%%%%%%%%%%%%%%%%%%%%%%%%%%%%%%%%%%%%%%%%%%%%%%%

\section{Validation and reproducibility}
\label{app:validation}

This appendix collects the checks the paper's results rest on, the two
occasions on which those checks failed, and the claims withdrawn as a result.

\begin{declbox}[frametitle={Two principles, each learned from a measured failure}]
\emph{An invariant tested at a point holds at that point.} A property verified
at one configuration carries no information about any other, and invariants
must be parametrised over the range a sweep will reach \emph{before} the sweep
rather than after.

\emph{A check that reports success without inspecting anything is worse than no
check}, because it converts ``not measured'' into ``measured and sound''. The
absence of a check is visible; a check that cannot fail is not.

Both are stated as principles because both were measured. The stock-flow
invariant of \cref{app:ownership} held only at the default configuration while
the test suite stayed green throughout, and two of this document's own build
detectors reported clean counts on a log that contained plenty
(\cref{app:repro}).
\end{declbox}

\subsection{An invariant that held only at the default}
\label{app:ownership}
\label{sec:ownership}

Firm ownership was assigned by cycling over \emph{households}: capitalist $i$
owned firm $i \bmod n_f$. That is a bijection \emph{only} at the default
configuration (100 households $\times$ 0.10 $=$ 10 capitalists, 10 firms). Off
the default it broke in both directions (\cref{tab:ownership}).

\begin{table}[htbp]
\centering
\small
\caption{The ownership defect, measured at seed 0 over 200 steps.}
\label{tab:ownership}
\begin{tabular}{lrrrr}
\toprule
$\pi_c$ & Capitalists & Firms with an owner & Money $t{=}0 \to t{=}200$
& $\sum$ net worth vs $K$\\
\midrule
0.05 & 5  & \textbf{5/10}  & $400.00 \to \mathbf{11.34}$ & $0.53\times$\\
0.08 & 8  & \textbf{8/10}  & $400.00 \to \mathbf{46.15}$ & $0.84\times$\\
0.10 & 10 & 10/10 & $400.00 \to 400.00$ & $1.05\times$\\
0.20 & 20 & 10/10 & $400.00 \to 400.00$ & $\mathbf{2.10\times}$\\
\bottomrule
\end{tabular}
\end{table}

Below the default, an ownerless firm's dividends and residual cash were paid
out inside a guard that tested for an owner, so they simply vanished --- a
direct violation of stock-flow consistency. Above it, the assignment overwrote
itself, leaving capitalists holding stale references whose capital was still
summed into net worth, inflating measured wealth inequality by a factor of two.

What makes the defect worth a subsection is \emph{when} it was found:
immediately before the global sensitivity analysis, which sweeps exactly the
parameter that exposed it. Undetected, \cref{sec:sa} would have reported
variance decompositions for a model that leaks money. The fix cycles over
\emph{firms}, so every firm has exactly one owner for any number of
capitalists; at the default the assignment is unchanged and a slice of the
committed panels reproduces exactly, so no previously reported result moves.

\subsection{The reproducibility criterion, and its two declared limits}
\label{app:bytecheck}

The nesting checks of \cref{sec:method} are verified
\emph{artifact-against-artifact} on committed panels. The original criterion --- deviation exactly $0.0$ --- proved \emph{not
reproducible over time}. Code that had recorded 7/7 exact reproductions later
deviated by 1 unit in the last place (ULP) on the same cells; eight hypotheses
were excluded by measurement and the cause remains unidentified. The magnitude
across 160 cells and 24 metrics is at most 2.1 ULP, does not amplify, and
produces \emph{zero} regime flips, so no economic conclusion moves. Re-running
the reference slice under the replacement gives 7/7 with a worst significant
drift of $0.00$ ULP --- it reproduces exactly \emph{today}, where the earlier
measurement found 2.1 ULP on the same code. That the drift is intermittent is
the strongest argument for retiring exact equality: a criterion that passes or
fails by the day carries no information about what changed in the model.

The replacement has two limbs that cannot compensate for each other: a declared
tolerance on \emph{levels} (8 ULP with an absolute floor, roughly four times
the measured envelope), and a check on \emph{regime} indicators --- viability,
which constraint binds, the sign of every resolvable metric --- at tolerance
exactly zero, since that is where drift would become scientific rather than
numerical. Two limits are declared in the source. A pure ULP tolerance is
unusable near zero, a metric legitimately resting at zero registering thousands
of ULP on a gap of $10^{-16}$; hence the absolute floor. And it must not be
applied to differenced or fitted quantities --- a chord, a slope, a curvature ---
because catastrophic cancellation makes their relative error arbitrary even
from stable inputs: the sensitivity analysis's slope deviates by 3410 ULP while
its own inputs deviate by 4. Those are checked by \emph{sign} instead.

\subsection{The bridging experiment in full}
\label{app:bridge}

\Cref{sec:estimator} reports the verdict; this is the design that produced it.
The experiment is a $2\times2$: method (chord versus OLS) crossed with
conditioning (all other parameters fixed at their defaults versus marginalised
over the sensitivity ranges), on a shared grid of runs. The verdict rule was
fixed \emph{before} execution: a sign flip across the method axis at fixed
parameters implicates the estimator; a flip across the conditioning axis at
constant method implicates conditioning; both flipping means both contribute
and the weights must be quantified; neither flipping is a finding to be
reported, not forced. The fixed-parameter arm cost no simulation, the committed
panels already carrying the canonical grid at twenty seeds, so only the
marginalised arm required new runs. \Cref{tab:bridge} reports the design.

\begin{table}[htbp]
\centering
\small
\caption{Bridging experiment, complete. The fixed arm holds every parameter at
the configuration in which \cref{sec:frontier} identified $\sigstar$
($c_0 = 1.0$, $\eta = 0$); the marginalised arm draws the other fifteen
parameters from the \cref{sec:sa} ranges on the same $\sigma$ grid. All rows
use the four-point $\rho$ support shared by both arms, so the comparison is not
confounded by the support.}
\label{tab:bridge}
\begin{tabular}{llrll}
\toprule
Cell & Parameters & Method & $\sigstar$ & 95\% CI\\
\midrule
1 & fixed        & chord & 0.447 & $[0.395,\,0.501]$\\
2 & fixed        & OLS   & 0.643 & $[0.596,\,0.701]$\\
3 & marginalised & chord & 0.938 & $[0.573,\,0.976]$\\
4 & marginalised & OLS   & \textbf{0.962} & $[0.723,\,0.985]$\\
\bottomrule
\end{tabular}
\end{table}

The mandatory anchoring control passes to fourteen significant figures: cell~2
evaluated on the full committed support returns
$\sigstar = 0.6540142777288407$ against the published $0.654$ with interval
$[0.616, 0.691]$. The bridge is therefore measuring the same object as
\cref{sec:frontier}, and any difference across the table is the treatment.

\paragraph{What the rule does not capture, reported anyway.} The rule asks only
whether the empirical band changes \emph{sides}, and by that test conditioning
does not fire: cells 2 and 4 both place the band below $\sigstar$. But
conditioning moves $\sigstar$ \emph{further} than the estimator does, and with
non-overlapping intervals. Reporting only the verdict would let a threshold
stand in for the measurement --- both displacements push the same way, and the
empirically supported band is profit-led more robustly than \cref{sec:frontier}
alone claimed.

\subsection{Withdrawn claims}
\label{app:retractions}

The project's declared rule is that every number entering a table must be
measured from a committed driver or labelled a design target
(\cref{sec:method}). Applying it retrospectively required withdrawing three
claims, recorded here rather than silently corrected.

\begin{enumerate}
  \item \textbf{A contradiction that did not exist.} An earlier summary
  reported an \emph{inverted} sign direction between the conditional and
  marginal analyses and left it open as a contradiction. There was none: ``the
  empirical band sits below $\sigstar$ and is therefore not wage-led'' and
  ``below $\sigma \approx 0.65$ wage-led outcomes are rare'' are the same
  statement. The error was in our prose, not in any measurement. The bridging
  experiment above was designed to explain the inversion and survived the
  discovery that half of what it set out to explain was a reading error --- its
  method axis still measured a real disagreement --- and is reported as it was
  fixed, not as it would have been written knowing the answer.
  \item \textbf{An inverted label.} A summary of the adaptive-expectations
  result claimed ``the wage-led result is robust to the expectation gain'',
  citing as evidence that the empirical band stays below $\sigstar$ --- which is
  the profit-led side. The measured $\lambda_e$-invariance is unaffected; the
  label attached to it was wrong, and \cref{sec:expectations} states it
  correctly.
  \item \textbf{A stale identity.} The relation $I/Y = \rho\alpha$ was carried
  forward from an earlier Cobb-Douglas core with no labour market and used to
  reason about anchoring after it had ceased to hold; measured $I/Y$ is cited
  in its place throughout (\cref{sec:calibration}). A related error mixed
  comparators --- business-sector $I/Y$ against whole-economy $K/Y$ --- and
  produced a $\delta$ that appeared anchored. Both are withdrawn.
\end{enumerate}

\subsection{Reproducibility}
\label{app:repro}

Every table in this paper comes from a committed output file with a committed
generator script; the correspondence is recorded in the repository. Runs pin
BLAS thread counts before importing numerical libraries, because
thread-dependent reduction order would otherwise break per-seed determinism,
and the execution environment --- library versions and the frozen sampling seed
--- is recorded alongside the sensitivity-analysis outputs.

The suite comprises 527 automated tests. Beyond conventional unit coverage they
assert the model's economic invariants directly: stock-flow consistency
parametrised over the ownership range rather than at the default alone
(\cref{app:ownership}), per-seed determinism, exact nesting at every limiting
parameter value, the workforce cap that restores diminishing returns
(\cref{sec:model}), and directional regression pins on basin selection. Fifteen
cover the repaired quantity of interest and the criterion of
\cref{app:bytecheck}, pinning the \emph{reason} for each declared constant
rather than the constant itself --- including, as an executable statement of the
defect this paper repairs, that a chord and a whole-support slope can carry
opposite signs on the same U-shaped response.

\paragraph{Two detectors that reported a reassuring zero.} The second measured
instance of this appendix's opening principle concerns the present document
rather than the model. The job that compiles the paper counts overfull and
underfull boxes in the \TeX{} log. Both counters matched the literal string
\verb|\hbox| through a layer of shell quoting that consumed the backslash, so
both matched nothing and both reported zero on logs containing many: a green
check that had inspected nothing, read as evidence of typographic cleanliness.
The counts quoted for this document were produced after the detectors were
repaired and verified against a synthetic log with known boxes.
